% =====================================================================
%  Artigo (autônomo): levantamento sistemático sobre coleta automatizada
%  de dados em Redes Sociais Digitais.
%
%  Compila com:  pdflatex survey && bibtex survey && pdflatex survey && pdflatex survey
%
%  Classe `article` padrão para compilar em qualquer instalação. Para
%  submissão, portar para o template do veículo (ex.: sbc-template) mantendo
%  o corpo das seções.
% =====================================================================
\documentclass[11pt,a4paper]{article}

\usepackage[utf8]{inputenc}
\usepackage[T1]{fontenc}
\usepackage[brazil]{babel}
\usepackage[a4paper,margin=2.5cm]{geometry}
\usepackage{graphicx}
\usepackage{booktabs}
\usepackage{amsmath}
\usepackage{enumitem}
\usepackage{abstract}
\usepackage[round]{natbib}
\usepackage[hidelinks]{hyperref}

\title{\textbf{Coleta, Filtragem e Volume em Análises de Redes Sociais\\
Digitais: um Levantamento Sistemático da Literatura}}
\author{Mariana Barreto \and Sérgio Lifschitz\thanks{Afiliação e contato a
completar.}\\[2pt]\small Departamento de Informática, PUC-Rio}
\date{\today}

\begin{document}
\maketitle

\begin{abstract}
\noindent
Análises computacionais sobre Redes Sociais Digitais (RSD) --- da Análise de
Redes Sociais (SNA) à modelagem de tópicos --- dependem de uma etapa de
\emph{coleta de dados} raramente problematizada. Sabe-se pouco, de forma
sistemática e em escala, sobre \emph{como} a literatura coleta, filtra e
persiste esses dados, e sobre o quanto essas escolhas reduzem o volume antes
mesmo da análise. Este artigo apresenta um levantamento sistemático e
reprodutível que, partindo do índice DBLP, cataloga 13.395 artigos e, após
extração estruturada, caracteriza a coleta em 1.718 que efetivamente coletam
dados de RSD. Os resultados revelam que a \emph{amostragem estatística} é rara
(\(\approx\)19\% dos artigos), enquanto a \emph{filtragem por critério} ---
que corta volume antes da análise --- é a norma (\(\approx\)80\%); que a
coleta via API domina (\(\approx\)70\%); e que a persistência de dados raramente usa
repositórios com DOI persistente. Um recorte sobre o \emph{volume} coletado
expõe ainda uma minoria de ``grandes coletores'', que reúnem de milhões a
bilhões de itens filtrando agressivamente --- quase sempre sem avaliar o
efeito dessas reduções sobre suas próprias conclusões. A extração dos dados de
coleta foi feita com modelos de linguagem validados contra anotação manual,
escolha metodológica que viabiliza a escala do levantamento e que detalhamos
para garantir reprodutibilidade.

\vspace{4pt}
\noindent\textbf{Palavras-chave:} redes sociais digitais; coleta de dados;
amostragem; análise de redes sociais; revisão sistemática; modelos de linguagem.
\end{abstract}

% ---------------------------------------------------------------------
\section{Introdução}

As Redes Sociais Digitais (RSD) tornaram-se um substrato central da
\textit{computational social science}~\citep{lazer2009computational}, e a
literatura que as estuda cresceu acompanhando o volume de dados que elas
geram. Toda análise sobre uma RSD --- detecção de comunidades, centralidade,
sentimento, modelagem de tópicos --- é precedida por uma etapa de \emph{coleta}.
Essa etapa, no entanto, é raramente o objeto de estudo: os trabalhos reportam
\emph{o que} coletaram, mas há pouco conhecimento sistemático sobre \emph{como}
o campo coleta, recorta e persiste seus dados, e sobre as consequências dessas
escolhas. Sabe-se, por exemplo, que amostrar uma rede pode distorcer suas
propriedades estruturais~\citep{leskovec2006sampling} e que a amostra entregue
por uma API pode divergir do todo~\citep{morstatter2013sample}; mas falta um
retrato em escala das práticas de coleta efetivamente adotadas.

Este artigo preenche essa lacuna com um \textbf{levantamento sistemático e
reprodutível}. A contribuição central é empírica: uma \emph{caracterização
quantitativa} de como 1.718 artigos coletam, filtram e persistem seus dados,
cujo achado principal é a raridade da amostragem estatística frente à
onipresença da filtragem por critério --- reduções de volume aplicadas antes
da análise e quase nunca avaliadas quanto a seu efeito sobre as conclusões.
Sustentam esse resultado dois instrumentos reprodutíveis: um \emph{pipeline}
que vai do índice bibliográfico à extração estruturada, com proveniência
rastreável; e um procedimento de extração em que um piloto confronta três LLMs
contra um gabarito anotado e retém os dois mais fiéis, mitigando alucinação por
triangulação.

% ---------------------------------------------------------------------
\section{Trabalhos Relacionados}

O estudo de RSD por bancos de dados e esquemas conceituais foi tratado em
trabalho anterior, que propôs uma especificação conceitual para múltiplas
plataformas e ilustrou análises típicas de rede e de
conteúdo~\citep{salgueiro2023,salgueiro2022dbmodels}. Aquele trabalho assume os
dados \emph{já coletados}; o presente levantamento volta-se à etapa anterior, a
da coleta, e a caracteriza empiricamente em escala. Quanto ao método, revisões
sistemáticas tradicionais dependem de triagem manual; aqui, a extração de campos
estruturados é delegada a LLMs com triangulação inter-modelo, abordagem ainda
pouco explorada para sintetizar práticas metodológicas de um corpus de milhares
de artigos. Em escala nacional, um levantamento manual dos eventos brasileiros
da área (Brasnam e SBBD) já havia inspecionado 100 artigos e apontado o
predomínio de uma única plataforma e o uso raro de SGBDs~\citep{heine2025database};
o presente trabalho generaliza e automatiza essa inspeção para um corpus
internacional de milhares de artigos.

% ---------------------------------------------------------------------
\section{Metodologia}

O levantamento é um \textit{pipeline} reprodutível: cada DOI catalogado tem
proveniência clara (a consulta que o retornou) e cada extração registra modelo,
\textit{prompt} e \textit{timestamp}.

\subsection{Coleta de DOIs e deduplicação}
A coleta primária usa a API de busca do DBLP~\citep{ley2002dblp}, percorrendo o
produto cartesiano de 6 plataformas (\texttt{instagram}, \texttt{twitter},
\texttt{facebook}, \texttt{youtube}, \texttt{reddit}, \texttt{tiktok}) por 17
anos (2010--2026). Foram 102 consultas e 17.812 \textit{hits}, deduplicados em
duas camadas --- por DOI e por título normalizado, preferindo a versão em acesso
aberto --- resultando em \textbf{13.365 candidatos únicos} da consulta ao DBLP
(10.230 em \textit{paywall}); somadas as inclusões posteriores, o catálogo
corrente contabiliza \textbf{13.395 artigos}.\footnote{Os números de coleta e
\textit{download} de PDFs referem-se ao recorte de maio de 2026; o catálogo
cresceu de forma marginal desde então.}

\subsection{\textit{Download} de PDFs}
Para cada DOI, um \textit{fetcher} agrega URLs candidatas de arXiv,
Unpaywall~\citep{piwowar2018unpaywall}, OpenAlex~\citep{priem2022openalex} e
Semantic Scholar, validando o PDF por \textit{Content-Type} e \textit{magic
bytes}. Foram baixados \textbf{2.106 PDFs} em acesso aberto (taxa de 67\%). Um
acesso institucional via \textit{proxy} recupera adicionalmente Springer e IEEE
(\(\approx\)43\% do \textit{paywall}); as demais editoras ficam atrás de
\textit{challenge} do Cloudflare, introduzindo um viés de cobertura por editora
(\S\ref{sec:limitacoes}).

\subsection{Extração estruturada via LLMs}
\label{sec:extracao}
No desenho da extração, cada PDF pode ser analisado por até \textbf{três LLMs
independentes} (Gemini 2.5 Flash, GPT-4o-mini e Claude Haiku
4.5)\footnote{Identificadores exatos: \texttt{gemini-2.5-flash},
\texttt{gpt-4o-mini} e \texttt{claude-haiku-4-5-20251001}.} com o mesmo
\textit{prompt} e vocabulário compartilhado, permitindo triangulação; a
\S\ref{sec:validacao} detalha quais modelos foram efetivamente retidos. O \textit{prompt} é um formulário de 11
campos com 5 regras no topo; as duas mais consequentes:

\begin{description}[leftmargin=1.4em,style=nextline]
  \item[Universo \textit{vs.} filtragem \textit{vs.} subconjunto] distinção
  estrita entre o que foi coletado, o critério que o recortou e o volume final.
  \item[Filtragem $\neq$ amostragem] amostragem estatística (aleatória,
  estratificada, \textit{bootstrap}\dots) difere de filtragem por critério
  (\textit{matching}, \textit{threshold}, palavra-chave, recorte temporal). Esta
  distinção é a chave dos resultados (\S\ref{sec:resultados}).
\end{description}

Os campos seguem o padrão \emph{``menciona? + detalhe''} (um \textit{booleano}
mais um campo-texto), que separa ``não menciona'' de ``falha de extração'' e
permite estatísticas limpas.

\subsection{Validação, piloto e seleção de modelos}
\label{sec:validacao}
Um artigo rico em metodologia foi anotado manualmente como \textbf{gabarito};
na validação, o Claude reproduziu todos os campos, o Gemini acertou a estrutura
com pequenas invenções, e o GPT-4o-mini inventou análises de SNA. Um piloto de
100 artigos confirmou o padrão: o GPT-4o-mini adicionou ``sentimento'' sem
evidência em 28 casos e ``centralidade'' em 26, e afirmou ``nenhum filtro'' em
35 casos onde os outros dois detectaram filtragem. Ele foi, então,
\textbf{descartado} do processamento em massa por qualidade insuficiente ---
ainda que fosse o mais barato dos três. Mantiveram-se o Gemini Flash e o Claude
Haiku, cuja concordância média ficou em torno de 85\% nos campos categóricos.

Como o Claude Haiku custa, por chamada, de três a cinco vezes mais que o Gemini
Flash,\footnote{Custo médio observado por artigo: \texttt{gemini-2.5-flash}
\(\approx\)R\$0,05; \texttt{claude-haiku-4-5-20251001} \(\approx\)R\$0,15 a
R\$0,28.} adotou-se um esquema de dois níveis: o \textbf{Gemini processou todo o
corpus} e serve de base para os resultados, enquanto o Claude Haiku foi aplicado
a um subconjunto (887 dos 2.139 extraídos) como validação e triangulação; nos
demais, os achados repousam sobre o Gemini, cuja fidelidade foi aferida contra o
gabarito e o piloto.

Por fim, a extração também identifica artigos que, embora retornados pela busca,
\emph{não realizam coleta de dados de RSD} para análise computacional --- estudos
puramente qualitativos, de entrevistas ou de posicionamento, entre outros. O
Gemini sinalizou \textbf{421 dos 2.139} nessa condição, os quais foram
\textbf{excluídos da caracterização}. Os resultados a seguir referem-se, portanto,
aos \textbf{1.718 artigos} que efetivamente coletam dados de RSD (o Claude Haiku
cobre 732 desses, \(\approx\)43\%, para triangulação).

% ---------------------------------------------------------------------
\section{Resultados}
\label{sec:resultados}

Esta seção traça o retrato quantitativo da coleta no corpus: quais plataformas
o campo estuda, como coleta e recorta os dados, que análises predominam e quão
raramente os dados são persistidos de forma reprodutível. Os números referem-se
aos \textbf{1.718 artigos} que efetivamente coletam dados de RSD
(\S\ref{sec:validacao}); as porcentagens das tabelas seguem a união inter-modelo
(``ao menos um modelo afirmou''), e o contraste amostragem \textit{vs.}\ filtragem
(adiante) é reportado sobre a base do Gemini, de cobertura total, com o Claude
Haiku (732 artigos) como validação.

\paragraph{Plataformas (Tab.~\ref{tab:plataformas}).} O Twitter domina (mais da
metade do corpus), reflexo de sua API acadêmica histórica, hoje descontinuada;
Reddit, YouTube e Facebook formam o segundo grupo.

\paragraph{Tipos de análise (Tab.~\ref{tab:analises}).} Predominam estatística
descritiva, sentimento e classificação; as análises de rede (comunidades,
centralidade, estrutura de rede, influência) somam fatia expressiva. O rótulo livre mais
citado foi ``análise de conteúdo'' (1.040 ocorrências), omitido da tabela por
ser um \emph{guarda-chuva} (engloba sentimento, tópico, enquadramento), e não um
método.

\paragraph{Estratégias de coleta (Tab.~\ref{tab:coleta}).} A coleta via API
domina (69,7\%), seguida de \textit{datasets} pré-existentes (45,8\%); o
\textit{scraping} aparece em apenas 19,1\%. As APIs mais citadas são as do
Twitter (várias formas), a YouTube Data API v3 e a do Reddit via PRAW, além da
extinta Pushshift~\citep{baumgartner2020pushshift}.

\begin{table}[t]
\centering
\begin{minipage}[t]{0.46\textwidth}\centering
\caption{Plataformas-alvo mais frequentes.}
\label{tab:plataformas}
\begin{tabular}{lr}
\toprule
Plataforma & Artigos \\ \midrule
Twitter & 1.095 \\ Reddit & 230 \\ YouTube & 204 \\
Facebook & 102 \\ TikTok & 88 \\ Instagram & 62 \\
\bottomrule
\end{tabular}
\end{minipage}\hfill
\begin{minipage}[t]{0.50\textwidth}\centering
\caption{Estratégia de coleta ($N=1.718$).}
\label{tab:coleta}
\begin{tabular}{lrr}
\toprule
Estratégia & Art. & \% \\ \midrule
API & 1.198 & 69,7 \\
\textit{Dataset} pré-existente & 787 & 45,8 \\
Outro método & 366 & 21,3 \\
\textit{Scraping} & 328 & 19,1 \\
\bottomrule
\end{tabular}
\end{minipage}
\end{table}

\begin{table}[t]
\centering
\caption{Tipos de análise (bucketizados), $N=1.718$.}
\label{tab:analises}
\begin{tabular}{lrr@{\hspace{2em}}lrr}
\toprule
Tipo & Art. & \% & Tipo & Art. & \% \\ \midrule
Estatística descritiva & 805 & 46,9 & Modelagem de tópico   & 394 & 22,9 \\
Sentimento             & 775 & 45,1 & Detecção de comunidades & 323 & 18,8 \\
Classificação          & 687 & 40,0 & Centralidade          & 304 & 17,7 \\
Temporal               & 629 & 36,6 & Influência            & 228 & 13,3 \\
Análise de rede        & 499 & 29,0 & Clusterização        & 148 & 8,6 \\
PLN (NER/\textit{embed.}) & 399 & 23,2 & Desinformação/\textit{bots} & 130 & 7,6 \\
\bottomrule
\end{tabular}
\end{table}

\paragraph{Amostragem é rara; filtragem é a norma.} Este é o achado central do
levantamento. Apenas \textbf{18,7\%} dos artigos reportam amostragem estatística,
enquanto \textbf{80,3\%} reportam filtragem por critério --- ou seja, a literatura
raramente ``amostra para inferir sobre um todo'' e quase sempre ``recorta o
objeto por critério'', tipicamente sem examinar se o recorte preserva as
conclusões. A direção do achado é robusta: os dois modelos concordam (Claude
Haiku: 21\% de amostragem e 91\% de filtragem sobre o subconjunto validado),
seguindo a distinção definida na \S\ref{sec:extracao}.

\paragraph{Persistência e reprodutibilidade.} A tecnologia de armazenamento é
mencionada em apenas \(\approx\)14\% dos artigos e, quando o é, predominam
formatos de arquivo simples (CSV, JSONL, Excel, Parquet) sobre SGBDs. Quanto à
publicação dos dados, o GitHub é o mais citado (\(\approx\)380 artigos), seguido
de disponibilização ``mediante solicitação'' (\(\approx\)140); repositórios com
DOI persistente (Zenodo, Figshare, OSF) aparecem em \(\approx\)140 --- cerca de
8\% do corpus, indicando baixa adesão aos princípios
FAIR~\citep{wilkinson2016fair}.

\subsection{Volume de coleta e os grandes coletores}
\label{sec:volume}

Quanto \emph{coletam} os trabalhos? A resposta é uma cauda longa: a coleta
típica fica na casa das dezenas de milhares de itens, mas uma minoria reúne de
milhões a bilhões --- e coletar muito, como se verá, não implica amostrar.
Medir isso exige o volume numérico do que foi coletado, disponível de forma
estruturada para \textbf{1.518 dos 1.718 artigos}\footnote{A quantidade coletada
é extraída pelos modelos como pares quantidade/unidade (campo
\texttt{collection\_items}) e reside na extração do Gemini, de cobertura total;
1.518 artigos têm ao menos uma quantidade legível. Tomamos como volume do artigo
a \emph{maior} quantidade entre seus itens de coleta, excluídas unidades que não
correspondem a itens coletados (\emph{tokens} de corpora de \textit{embeddings},
contadores de visualização, palavras e sentenças de \textit{snippets}), para não
inflar o volume com métricas alheias à coleta.}. Como as unidades variam
(\textit{tweets}, vídeos, comentários, usuários\dots), esse indicador da maior
dimensão coletada é robusto à mistura de unidades.

\paragraph{Uma cauda longa de quatro ordens de grandeza (Fig.~\ref{fig:volume}).}
O volume mediano é de \(\approx\)273 mil itens, mas a distribuição se estende por
mais de sete ordens de grandeza: 19\% dos artigos coletam menos de 10 mil itens,
enquanto 40\% ultrapassam 1 milhão e \textbf{339 artigos (22\%) coletam ao menos
10 milhões}. No extremo estão coletas quase exaustivas ou de acesso
\emph{privilegiado}: \(\approx\)118 bilhões de \textit{tweets} via
\textit{decahose} (amostra de 10\%) da API do Twitter ao longo de uma década,
\(\approx\)30,9 bilhões de arestas de um grafo interno do Twitter, ou
\(\approx\)13 bilhões de comentários do Reddit desde 2005. Esses extremos
refletem condição de acesso (infraestrutura interna, \emph{firehose}, dumps
completos), não coleta típica via API pública --- um lembrete de que o topo da
cauda reflete condição de acesso, não técnica de coleta reproduzível.

\begin{figure}[t]
\centering
\includegraphics[width=0.82\textwidth]{fig_volume.pdf}
\caption{Distribuição do volume coletado por ordem de grandeza (maior
quantidade por artigo), \(N=1.518\) artigos com volume legível. Em destaque, os
339 \emph{grandes coletores} (\(\geq\)10 milhões de itens).}
\label{fig:volume}
\end{figure}

\paragraph{Coletar muito não implica amostrar (Tab.~\ref{tab:grandes}).} Entre os
339 grandes coletores, a amostragem estatística dentro do universo coletado
segue \emph{minoritária}: apenas \textbf{80 (24\%)} amostram, coerente com o
achado geral (\S\ref{sec:resultados}). Quando amostram, a fração efetivamente
analisada é diminuta: um estudo coletou \(\approx\)2,8 bilhões de \textit{tweets}
e analisou 28 milhões (\(\approx\)1\%); outro coletou 500 milhões e reteve
28 mil (\(\approx\)0,01\%). Ou seja, mesmo na coleta em escala massiva, o objeto
de análise costuma ser um recorte definido por critério --- \textbf{297 dos 339
(88\%)} relatam filtragem ---, sem verificação de que o recorte preserva as
conclusões.

\paragraph{Escala não traz reprodutibilidade.} A infraestrutura de dados não
acompanha o volume: apenas \textbf{24 dos 339} (7\%) mencionam qualquer
tecnologia de armazenamento, e \textbf{149 (44\%)} não dão nenhuma indicação de
persistência dos dados coletados --- apesar de terem reunido dezenas de milhões
a bilhões de itens. A publicação em repositório com DOI persistente aparece em
apenas 36 casos (11\%); somada ao GitHub, chega a 100 (30\%). A precariedade de
infraestrutura observada no corpus (\S\ref{sec:resultados}) persiste, portanto,
também no extremo superior de volume.

\begin{table}[t]
\centering
\caption{Comportamento dos 339 grandes coletores (\(\geq\)10 milhões de itens).}
\label{tab:grandes}
\begin{tabular}{lrr}
\toprule
& Artigos & \% \\ \midrule
Amostram dentro do universo coletado & 80  & 24 \\
Filtram por critério                 & 297 & 88 \\
Mencionam tecnologia de BD           & 24  & 7 \\ \midrule
\textit{Persistência:} DOI/repositório   & 36  & 11 \\
\textit{Persistência:} GitHub            & 64  & 19 \\
\textit{Persistência:} sob solicitação   & 21  & 6 \\
\textit{Persistência:} outra menção      & 69  & 20 \\
\textit{Persistência:} nenhuma indicação & 149 & 44 \\
\bottomrule
\end{tabular}
\end{table}

% ---------------------------------------------------------------------
\section{Discussão}

O retrato é o de um campo que \emph{recorta muito e justifica pouco}: a coleta é
quase sempre um subconjunto definido por critério (palavra-chave, janela,
\textit{top-k} de API, comunidade), sem verificação de que esse recorte é
inócuo para as conclusões. Combinado ao que se sabe sobre o efeito da amostragem
em grafos~\citep{leskovec2006sampling} e à divergência entre amostra de API e
todo~\citep{morstatter2013sample}, isso sugere um risco sistemático de
conclusões dependentes da estratégia de coleta --- hipótese que motiva
investigação experimental dedicada (\S\ref{sec:conclusao}). Esse retrato
converge com o de um levantamento manual precursor, restrito aos eventos
brasileiros da área (83 artigos do Brasnam e 17 do SBBD), que já observara o
predomínio do Twitter/X, análises conduzidas sobre amostras do todo e uso raro
de SGBDs~\citep{heine2025database}; o presente trabalho estende essa observação
a milhares de artigos e à escala internacional. No plano de
reprodutibilidade, a baixa adesão a repositórios persistentes e o predomínio de
arquivos planos sobre SGBDs indicam que a infraestrutura de dados do campo é, em
geral, \textit{ad hoc}.

\subsection{Limitações}
\label{sec:limitacoes}
\begin{enumerate}[leftmargin=1.4em]
  \item \textbf{Cobertura do DBLP}: o universo limita-se ao que o DBLP indexa;
  revistas de comunicação/sociologia, fora do escopo de ciência da computação,
  podem ficar de fora.
  \item \textbf{Viés de cobertura por editora}: dos 10.230 artigos em
  \textit{paywall}, o acesso institucional via \textit{proxy} recupera apenas
  Springer e IEEE (\(\approx\)43\% do \textit{paywall}); ACM, Wiley, SAGE,
  Taylor \& Francis e Emerald ficam atrás do \textit{challenge} do Cloudflare, e
  a Elsevier atrás de anti-\textit{bot} próprio, permanecendo inacessíveis. Como
  a extração recai sobre os PDFs obtidos, essas editoras ficam sub-representadas.
  \item \textbf{Subconjunto extraído}: os 1.718 artigos caracterizados (dos 2.139
  com PDF acessível e extração estruturada, após excluir 421 que não coletam dados
  de RSD) são os de PDF acessível (acesso aberto mais Springer/IEEE via
  \textit{proxy}), não uma amostra aleatória dos 13.395 catalogados; os achados
  descrevem essa fração acessível, e o viés de editora acima se propaga aos
  resultados.
  \item \textbf{Não determinismo e alucinação dos LLMs}: mitigados por
  triangulação de modelos, validação contra gabarito manual e versionamento de
  cada extração. A concordância inter-modelo (\(\approx\)85\% nos campos
  categóricos) delimita a confiança nos achados; o campo mais subjetivo --- o
  tipo de análise --- é também o de maior divergência entre modelos.
\end{enumerate}

% ---------------------------------------------------------------------
\section{Conclusão}
\label{sec:conclusao}

Este levantamento caracterizou, em escala (13.395 artigos catalogados, 1.718 com
coleta de RSD caracterizada) e de forma reprodutível, como a literatura de RSD coleta
seus dados, com o achado central de que a amostragem estatística é rara e a
filtragem por critério é a norma --- quase sempre sem verificação de impacto.
Esse resultado abre uma agenda experimental clara:
medir, por tipo de análise, o efeito de coletar mais do que os trabalhos
coletaram, testando se as conclusões publicadas se sustentam sob coletas mais
amplas --- agenda desenvolvida em trabalho subsequente.

\bibliographystyle{plainnat}
\bibliography{referencias}

\end{document}
